%==============================================================================
% Sjabloon onderzoeksvoorstel bachproef
%==============================================================================
% Gebaseerd op document class `hogent-article'
% zie <https://github.com/HoGentTIN/latex-hogent-article>

% Voor een voorstel in het Engels: voeg de documentclass-optie [english] toe.
% Let op: kan enkel na toestemming van de bachelorproefcoördinator!
\documentclass{hogent-article}

% Packages voor figuren
\usepackage{graphicx}
\usepackage{float}

% Invoegen bibliografiebestand
\addbibresource{voorstel.bib}

% Informatie over de opleiding, het vak en soort opdracht
\studyprogramme{Professionele bachelor toegepaste informatica}
\course{Bachelorproef}
\assignmenttype{Onderzoeksvoorstel}
% Voor een voorstel in het Engels, haal de volgende 3 regels uit commentaar
% \studyprogramme{Bachelor of applied information technology}
% \course{Bachelor thesis}
% \assignmenttype{Research proposal}

\academicyear{2025-2026} % TODO: pas het academiejaar aan

% TODO: Werktitel
\title{AI-gestuurd Dependency Dashboard voor Automatische Analyse en Prioritering van Software-updates}

% TODO: Studentnaam en emailadres invullen
\author{Jelle Vandriessche}
\email{jelle.vandriessche@student.hogent.be}

% TODO: Medestudent
% Gaat het om een bachelorproef in samenwerking met een student in een andere
% opleiding? Geef dan de naam en emailadres hier
% \author{Yasmine Alaoui (naam opleiding)}
% \email{yasmine.alaoui@student.hogent.be}

% TODO: Geef de co-promotor op
\supervisor[Co-promotor]{B. Pattyn (Springbok Agency, \href{mailto:bert.pattyn@springbokagency.com}{bert.pattyn@springbokagency.com})}
% Github Repository link
\projectrepo{https://github.com/jellev00/hogent-bachproef-paper-jelle-vandriessche}

% Binnen welke specialisatierichting uit 3TI situeert dit onderzoek zich?
% Kies uit deze lijst:
%
% - Mobile \& Enterprise development
% - AI \& Data Engineering
% - Functional \& Business Analysis
% - System \& Network Administrator
% - Mainframe Expert
% - Als het onderzoek niet past binnen een van deze domeinen specifieer je deze
%   zelf
%
\specialisation{AI \& Software Engineering}
\keywords{AI-agent, dependency management, software maintenance, technische schuld, geautomatiseerde updates}

\begin{document}

\begin{abstract}
  Deze bachelorproef ontwikkelt een aparte webapplicatie (dashboard) met een from-scratch ontworpen AI-agent om dependency-updates automatisch te detecteren, analyseren en prioriteren voor het kernteam dat platform- en dependencybeheer opneemt binnen een bedrijf. Vanuit de vastgestelde problematiek — veel manueel werk bij het interpreteren van changelogs, ophopende technische schuld en beperkte context bij updates — controleert de applicatie GitHub-repositories op gebruikte dependencies, haalt versie-informatie op uit publieke registries (zoals npm) en laat de AI-agent changelogs en release notes samenvatten met indicatie van kriticiteit (security, bugfix, feature) en risico-impact.

  Via een systematische aanpak met requirements-analyse, architectuurontwerp, proof-of-concept implementatie en gerichte vergelijking van AI-aanpakken wordt onderzocht hoe deze oplossing de manuele beoordelingslast kan verminderen en het overzicht op technische schuld kan verbeteren. De evaluatie met het kernteam focust op bruikbaarheid van het dashboard, duidelijkheid van AI-samenvattingen en ondersteuning bij prioritering. Verwacht wordt dat de aanpak een efficiënter, transparanter dependencybeheerproces oplevert dat herhaalbaar is en kan worden opgeschaald naar andere teams binnen de organisatie.
\end{abstract}

\tableofcontents

% De hoofdtekst van het voorstel zit in een apart bestand, zodat het makkelijk
% kan opgenomen worden in de bijlagen van de bachelorproef zelf.
\input{voorstel-inhoud}

\printbibliography[heading=bibintoc]

\end{document}